\documentclass[12pt]{article}
\usepackage[spanish]{babel}
\usepackage[utf8]{inputenc}
\usepackage[T1]{fontenc}
\usepackage{graphicx}
\usepackage{geometry}
\usepackage{setspace}
\usepackage[hidelinks]{hyperref}
\usepackage{float}
\usepackage{titlesec}
\usepackage{listings}
\usepackage{xcolor}
\usepackage{enumitem}
\usepackage{amsmath}

\geometry{a4paper, margin=2.5cm}

% --- Configuracion para codigo ---
\definecolor{codegreen}{rgb}{0,0.6,0}
\definecolor{codegray}{rgb}{0.5,0.5,0.5}
\definecolor{backcolour}{rgb}{0.95,0.95,0.92}

\lstset{
    backgroundcolor=\color{backcolour},   
    commentstyle=\color{codegreen},
    keywordstyle=\color{blue},
    numberstyle=\tiny\color{codegray},
    stringstyle=\color{orange},
    basicstyle=\ttfamily\footnotesize,
    breaklines=true,
    captionpos=b,
    numbers=left,
}

\begin{document}
% --- PORTADA ---
\begin{titlepage}
\begin{center}
    {\large \textbf{UNIVERSIDAD TÉCNICA ESTATAL DE QUEVEDO}} \\[0.5cm]
    {\large \textbf{FACULTAD DE CIENCIAS DE LA COMPUTACIÓN Y DISEÑO DIGITAL}} \\[0.5cm]
    {\large \textbf{CARRERA DE SOFTWARE}} \\[2cm]
    
    \textbf{TEMA:}\\[0.5cm]
    {\Large \textbf{Evaluación Aplicacion Web}} \\[1cm]
    
    \textbf{ALUMNO:}\\[0.5cm]
    \textbf{CRUZ PEREZ JUSTYN KEITH} \\[1cm]

    \textbf{REPOSITORIO DE GITHUB:}\\[0.3cm]
    \href{https://github.com/keithdrox/inventario-mercado-cruz}{\texttt{https://github.com/keithdrox/inventario-mercado-cruz}} \\[1cm]
    
    \textbf{DOCENTE:}\\[0.5cm]
    ING. GUERRERO ULLOA GLEISTON CICERON\\[1cm]
    
    \textbf{MATERIA:}\\[0.5cm]
    APLICACIONES WEB - A - [5to Nivel] \\[1cm]
    
    \vfill
    Quevedo - Ecuador \\
    2026
\end{center}
\end{titlepage}

\newpage
\tableofcontents
\newpage

\section{Introducción}
El presente informe detalla el diseño e implementación del módulo de gestión de inventario para el Mercado Central del Municipio de Quevedo. El sistema se expone mediante una API REST desarrollada con Spring Boot 3 y Java 21, siguiendo un esquema estructurado de controlador, servicio y repositorio. Se implementan funcionalidades clave como validación de datos, borrado lógico, seguridad basada en JSON Web Tokens (JWT) \cite{jwt}, y un mecanismo de caché utilizando Redis. Todo el entorno está orquestado mediante Docker Compose para garantizar la reproducibilidad y facilidad de ejecución.

\section{Arquitectura y Modelo de Datos}
La solución se basa en una arquitectura por capas bien definida, separando las responsabilidades de cada componente para mantener un diseño limpio y mantenible \cite{cleanarch}. 

\subsection{Capas de la Aplicación}
\begin{itemize}
    \item \textbf{Controladores:} Exponen los endpoints REST y gestionan las peticiones HTTP. 
    \item \textbf{Servicios:} Contienen la lógica de negocio, reglas de validación adicionales y orquestan la interacción con los repositorios y la caché.
    \item \textbf{Repositorios:} Punto de acceso a los datos, extendiendo de \texttt{JpaRepository}.
    \item \textbf{DTOs y Entidades:} Separación clara entre el modelo interno y los objetos expuestos al cliente (\texttt{ProductoRequest}, \texttt{ProductoResponse}, \texttt{ApiResponse}).
\end{itemize}

\subsection{Modelo de Datos}
El esquema se crea dinámicamente con DDL a través del contenedor PostgreSQL:

\begin{lstlisting}[language=SQL, caption=Esquema de base de datos]
CREATE TABLE usuarios (
    id BIGSERIAL PRIMARY KEY,
    username VARCHAR(50) NOT NULL UNIQUE,
    password VARCHAR(255) NOT NULL,
    rol VARCHAR(20) NOT NULL
);

CREATE TABLE productos (
    id BIGSERIAL PRIMARY KEY,
    nombre VARCHAR(100) NOT NULL,
    categoria VARCHAR(50) NOT NULL,
    stock INTEGER NOT NULL CHECK (stock >= 0),
    precio DECIMAL(10,2) NOT NULL CHECK (precio >= 0.01),
    activo BOOLEAN NOT NULL DEFAULT TRUE,
    creado_en TIMESTAMPTZ NOT NULL DEFAULT now()
);
\end{lstlisting}

\section{Implementación de Requisitos}

\subsection{Listado Paginado}
El endpoint \texttt{GET /api/v1/productos} devuelve una lista paginada de productos activos. Spring Data JPA facilita la paginación inyectando \texttt{Pageable}. La respuesta se envuelve en una estructura estándar JSON que incluye metadatos (\texttt{totalElements}, \texttt{totalPages}).

\subsection{Creación Validada}
Para la creación de productos mediante \texttt{POST}, se implementa validación a nivel de servidor utilizando \texttt{jakarta.validation} como \texttt{@NotBlank}, \texttt{@Min}, y \texttt{@DecimalMin}. 
Un \texttt{@RestControllerAdvice} intercepta excepciones y construye una respuesta estructurada estilo Problem Details \cite{problemdetails}, devolviendo estado HTTP 400.

\subsection{Eliminación Lógica (Soft Delete)}
Al eliminar mediante \texttt{DELETE}, se realiza un borrado lógico actualizando el campo \texttt{activo} a \texttt{false}. Si el producto no existe, lanza una \texttt{ResourceNotFoundException} y retorna 404.

\subsection{Cache-Aside con Redis}
Se integra Redis utilizando el patrón cache-aside:
\begin{itemize}
    \item \textbf{Lectura:} Anotación \texttt{@Cacheable} en el listado paginado. 
    \item \textbf{Invalidación:} Anotación \texttt{@CacheEvict(allEntries=true)} en creación y eliminación para limpiar la caché.
\end{itemize}

\subsection{Seguridad con JSON Web Tokens (JWT)}
Los endpoints están protegidos por Spring Security.
\begin{itemize}
    \item \texttt{GET} exige al menos el rol \texttt{ROLE\_USER}.
    \item \texttt{POST} y \texttt{DELETE} exigen \texttt{ROLE\_ADMIN}.
\end{itemize}
Ante una petición sin token, se devuelve HTTP 401. Si el rol es insuficiente, HTTP 403.

\section{Conclusión}
La API desarrollada cumple con los requerimientos técnicos y estándares modernos de desarrollo, garantizando reproducibilidad y solidez gracias a la contenedorización, seguridad JWT y caché con Redis.

\bibliographystyle{plain}
\bibliography{referencias}

\end{document}